\chapter{Solver command-line reference}
\label{chap:solverref}

\begin{Snippet}
lsasolver, Ver. 0.5.0 9/22/15, Run 2015/10/08 at 11:32:53
Usage: lsasolver [option] ...
 Program lsasolver will read an input file and ...
 Input is on the command line, or of the same format in a file (see --file below);
 lines in that file which begin with '#' are ignored. Accepted options are 
 shown below, followed by a description, with default value, if any, in ().

# File I/O:
 --file <name>       Name of file containing more options [#-EOL = comment]
 --out <name>        Name of output file ()
 --outpath <path>    Path for output file ()
 --input <name>      Name of input file(s) ()
 --inpath <path>     Path for input file ()

# Program control:
 --apquit            Quit after computing a priori positions (don't)

# Algorithm [*overwrite DAT file input; default if no DAT input]:
 --niter <n>         Limit on the number of iterations (default 15) [*] (-1)
 --conv <frac>       Convergence criterion (default 1.0e-08) [*] (0.00)
 --2D                Solve a true 2D problem (ignore Z) (don't)
 --noAPV             Leave covariance in relative units (default F) [*] (don't)
 --APV               Scale covariance to physical units (default T) [*] (don't)
 --alpha <prob>      Significance level of Chi-squared test [0<prob<1] (0.050)
 --SOA               Express angle equations in seconds-of-arc (don't)
 --allowCOM          Allow a priori computation to return center-of-mass solution
 --statsAll          Output chi squared and data snooping at every iteration
 --calcExtRelVect    Calculate external reliability values

# Geoid:
 --geoidfile <name>  Name of geoid file ()
 --geoidpath <path>  Path for geoid file ()
 --interpolation <method> Geoid interpolation method [bicubic,bilinear] (bicubic)

# Output:
 --eqnout <file>     Output observation equations to this file ()
 --bin <file>        Output results for GUI to file in CSV format ()
 --csv <file>        Output final positions and errors to CSV format file ()
 --extr <"str">      Extraction string (EXTR tag TYP label[s]) [repeat] ()
 --westLon           Output west longitude (don't)
 --datout <file>     Output a complete DAT file after adjustment ()
 --progress          Output to stdout summary information, incl. at each iteration

# Help:
 --validate          Read input and test its validity, then quit (don't)
 --verbose           Print extended output information (don't)
 --debug             Print debug output at level 0 [debug<n> for level n=1-7] (-1)
 --timing            Print timing information (don't)
 --help              Print this and quit (don't)
\end{Snippet}
